\section*{Chapitre \ref{ch:g12:meiosis-genetic-shuffling} --- La méiose et le brassage génétique}

\begin{solution}{exo:g12:meiosis-genetic-shuffling:1}
\omterm{def:g12:meiosis-genetic-shuffling:meiosis}{Diploïde}~: deux exemplaires de chaque \omterm{def:g11:cell-cycle-mitosis:chromosome}{chromosome} ($2n$)~; \omterm{def:g12:meiosis-genetic-shuffling:meiosis}{haploïde}~: un
seul ($n$). La \omterm{def:g12:meiosis-genetic-shuffling:meiosis}{méiose} divise par deux le nombre de \omterm{def:g10:universal-dna:information}{chromosomes}, de $2n$
à $n$, et ramène l'\omterm{def:g10:universal-dna:information}{ADN} de $2q$ (après réplication) à $q/2$ par gamète.
\end{solution}

\begin{solution}{exo:g12:meiosis-genetic-shuffling:2}
La \omterm{def:g12:meiosis-genetic-shuffling:meiosis}{méiose} I sépare les chromosomes homologues de chaque paire~; la
\omterm{def:g12:meiosis-genetic-shuffling:meiosis}{méiose} II sépare les deux \omterm{def:g11:cell-cycle-mitosis:chromosome}{chromatides} de chaque \omterm{def:g11:cell-cycle-mitosis:chromosome}{chromosome}. La première
est la division réductionnelle.
\end{solution}

\begin{solution}{exo:g12:meiosis-genetic-shuffling:3}
Après la \omterm{def:g12:meiosis-genetic-shuffling:meiosis}{méiose} I~: 4 \omterm{def:g10:universal-dna:information}{chromosomes} (à deux \omterm{def:g11:cell-cycle-mitosis:chromosome}{chromatides} chacun), \omterm{def:g10:universal-dna:information}{ADN} $q$.
Après la \omterm{def:g12:meiosis-genetic-shuffling:meiosis}{méiose} II~: 4 \omterm{def:g10:universal-dna:information}{chromosomes} (à une \omterm{def:g11:cell-cycle-mitosis:chromosome}{chromatide} chacun), \omterm{def:g10:universal-dna:information}{ADN} $q/2$.
\end{solution}

\begin{solution}{exo:g12:meiosis-genetic-shuffling:4}
$2^4 = 16$.
\end{solution}

\begin{solution}{exo:g12:meiosis-genetic-shuffling:5}
Un croisement avec un partenaire ne portant que des \omterm{def:g10:universal-dna:gene}{allèles} \omterm{def:g11:genes-and-disease:genetic}{récessifs},
de sorte que chaque descendant montre le gamète du parent testé. La
proportion de recombinés révèle si deux \omterm{def:g10:universal-dna:gene}{gènes} sont sur le même
\omterm{def:g11:cell-cycle-mitosis:chromosome}{chromosome} et, le cas échéant, à quelle distance l'un de l'autre.
\end{solution}

\begin{solution}{exo:g12:meiosis-genetic-shuffling:6}
\omterm{def:g12:meiosis-genetic-shuffling:meiosis}{Méiose} I~: $2q$~; entre les divisions~: $q$~; gamète~: $q/2$. Sans
phase S intercalée, la seconde division divise de nouveau l'\omterm{def:g10:universal-dna:information}{ADN} par
deux, ce qui amène le gamète à la moitié du contenu d'une \omterm{def:g10:cells-common-unit:cell}{cellule}
\omterm{def:g12:meiosis-genetic-shuffling:meiosis}{diploïde}.
\end{solution}

\begin{solution}{exo:g12:meiosis-genetic-shuffling:7}
Quatre classes égales~: \omterm{def:g10:universal-dna:gene}{gènes} non liés, sur des \omterm{def:g10:universal-dna:information}{chromosomes} différents.
Gamètes $AB$, $Ab$, $aB$, $ab$, un quart chacun.
\end{solution}

\begin{solution}{exo:g12:meiosis-genetic-shuffling:8}
Liés~: deux grandes classes ($AB$, $ab$, parentales) et deux petites
($Ab$, $aB$, recombinées). Recombinés $85/1000 = 8.5\%$.
\end{solution}

\begin{solution}{exo:g12:meiosis-genetic-shuffling:9}
Seulement $AB$ et $ab$~: un échange au-dessus des deux \omterm{def:g10:universal-dna:gene}{gènes} intervertit
des segments qui n'en portent aucun, si bien que les combinaisons de
$A/a$ et de $B/b$ sont inchangées. Une recombinaison entre deux \omterm{def:g10:universal-dna:gene}{gènes}
exige un échange entre eux.
\end{solution}

\begin{solution}{exo:g12:meiosis-genetic-shuffling:10}
Environ 0.8, 2.9 et 16 pour 1000~: un facteur 20 entre 25 et 42 ans.
\end{solution}

\begin{solution}{exo:g12:meiosis-genetic-shuffling:11}
En \omterm{def:g12:meiosis-genetic-shuffling:meiosis}{méiose} I, les deux homologues gagnent une même \omterm{def:g10:cells-common-unit:cell}{cellule}~: deux
gamètes à deux exemplaires et deux sans aucun. En \omterm{def:g12:meiosis-genetic-shuffling:meiosis}{méiose} II, les deux
\omterm{def:g11:cell-cycle-mitosis:chromosome}{chromatides} d'un \omterm{def:g11:cell-cycle-mitosis:chromosome}{chromosome} gagnent une même \omterm{def:g10:cells-common-unit:cell}{cellule}~: un gamète à deux
exemplaires, un sans aucun, et deux normaux.
\end{solution}

\begin{solution}{exo:g12:meiosis-genetic-shuffling:12}
La recombinaison est proportionnelle à la distance~: 2~\% signifie que
les \omterm{def:g10:universal-dna:gene}{gènes} sont très proches, si bien qu'un échange tombe rarement entre
eux~; 48~\% signifie qu'ils sont éloignés, avec un échange entre eux à
presque chaque \omterm{def:g12:meiosis-genetic-shuffling:meiosis}{méiose}. Comme un échange donne 50~\% de gamètes
recombinés, des \omterm{def:g10:universal-dna:gene}{gènes} liés mais éloignés se recombinent aussi
librement que des \omterm{def:g10:universal-dna:gene}{gènes} indépendants.
\end{solution}

\begin{solution}{exo:g12:meiosis-genetic-shuffling:13}
$AB$, $Ab$, $aB$, $ab$, un quart chacun, par brassage indépendant des
deux paires. Le \omterm{prop:g12:meiosis-genetic-shuffling:crossingover}{crossing-over} n'y change rien~: il recombine les
\omterm{def:g10:universal-dna:gene}{allèles} le long d'un même \omterm{def:g11:cell-cycle-mitosis:chromosome}{chromosome}, et ceux-ci sont sur des
\omterm{def:g10:universal-dna:information}{chromosomes} différents.
\end{solution}

\begin{solution}{exo:g12:meiosis-genetic-shuffling:14}
Chaque parent transmet à chaque enfant l'un de ses deux \omterm{def:g10:universal-dna:gene}{allèles} à
chaque \omterm{def:g10:universal-dna:gene}{gène}, tiré au hasard~; à un \omterm{def:g10:universal-dna:gene}{gène} donné, deux frères ou sœurs ont
reçu le même \omterm{def:g10:universal-dna:gene}{allèle} d'un parent avec une probabilité de $1/2$. Sur de
nombreux \omterm{def:g10:universal-dna:gene}{gènes}, la fraction partagée avoisine $1/2$, mais pour une
paire donnée elle fluctue autour de cette valeur.
\end{solution}

\begin{solution}{exo:g12:meiosis-genetic-shuffling:15}
Des clones de la plante mère, génétiquement identiques à elle et entre
eux. La plante y gagne une multiplication rapide et sûre d'un \omterm{def:g11:enzymes-and-phenotype:phenotype}{génotype}
qui fonctionne ici et maintenant~; elle y perd la variété qui
permettrait à certains descendants de survivre à un changement
d'environnement ou à un nouveau parasite.
\end{solution}

\begin{solution}{pb:g12:meiosis-genetic-shuffling:1}
\textbf{1.} Femelle $b^+b$, $vg^+vg$, avec $b^+vg^+$ sur un \omterm{def:g11:cell-cycle-mitosis:chromosome}{chromosome}
et $b\,vg$ sur l'autre (venus de ses deux parents de lignée pure).
Mâle $bb$, $vg\,vg$.

\textbf{2.} Uniquement des gamètes $b\,vg$~: le mâle n'apporte que des
\omterm{def:g10:universal-dna:gene}{allèles} \omterm{def:g11:genes-and-disease:genetic}{récessifs}, si bien que le \omterm{def:g11:enzymes-and-phenotype:phenotype}{phénotype} de chaque descendant révèle
le gamète de la femelle.

\textbf{3.} $b^+vg^+$ (gris-normal), $b\,vg$ (noir-vestigial),
$b^+vg$ (gris-vestigial), $b\,vg^+$ (noir-normal).

\textbf{4.} Gris-normal 42~\%, noir-vestigial 41~\%, gris-vestigial
9~\%, noir-normal 8~\%.

\textbf{5.} Sur le même \omterm{def:g11:cell-cycle-mitosis:chromosome}{chromosome}~: deux grandes classes et deux
petites au lieu de $1:1:1:1$.

\textbf{6.} Gris-normal et noir-vestigial~: les combinaisons que
portaient ses deux homologues, héritées intactes de ses parents de
lignée pure.

\textbf{7.} Gris-vestigial et noir-normal, nés d'un \omterm{prop:g12:meiosis-genetic-shuffling:crossingover}{crossing-over} entre
les deux \omterm{def:g10:universal-dna:gene}{gènes} en prophase I.

\textbf{8.} $(206 + 185)/2300 \approx 17\%$. Un \omterm{prop:g12:meiosis-genetic-shuffling:crossingover}{crossing-over} entre les
\omterm{def:g10:universal-dna:gene}{gènes} rend recombinées deux des quatre \omterm{def:g11:cell-cycle-mitosis:chromosome}{chromatides}, si bien que 17~\%
de gamètes recombinés signifient un échange dans environ 34~\% des
\omterm{def:g12:meiosis-genetic-shuffling:meiosis}{méioses}.

\textbf{9.} Un échange unique produit ensemble les deux \omterm{def:g11:cell-cycle-mitosis:chromosome}{chromatides}
recombinées, une de chaque sorte~; et les deux homologues gagnent les
gamètes à égalité, si bien que les deux classes parentales
s'équilibrent aussi.

\textbf{10.} Les classes parentales sont désormais gris-vestigial et
noir-normal, environ 41.5~\% chacune~; les recombinées gris-normal et
noir-vestigial, environ 8.5~\% chacune.

\textbf{11.} Oui~: une recombinaison bien inférieure à 50~\% avec les deux.

\textbf{12.} $b$ --- $cn$ --- $vg$~: 9 et 8 font bien 17.

\textbf{13.} La probabilité d'un échange dans un intervalle est
proportionnelle à sa longueur, si bien que les fréquences
d'intervalles adjacents s'additionnent~: la fréquence de recombinaison
mesure la distance le long du \omterm{def:g11:cell-cycle-mitosis:chromosome}{chromosome}.

\textbf{14.} Un échange donne deux \omterm{def:g11:cell-cycle-mitosis:chromosome}{chromatides} recombinées sur quatre~:
50~\% au plus~; des échanges supplémentaires rebrassent mais n'élèvent
jamais la fraction recombinée au-dessus de la moitié.

\textbf{15.} Il est sur un autre \omterm{def:g11:cell-cycle-mitosis:chromosome}{chromosome} ou très loin sur le même~;
les croisements ne permettent pas de trancher.

\textbf{16.} $2^4 = 16$ gamètes~; $16 \times 16 = 256$ combinaisons.

\textbf{17.} Chaque échange produit des \omterm{def:g10:universal-dna:information}{chromosomes} qui n'avaient
jamais existé, à une position qui varie d'une \omterm{def:g12:meiosis-genetic-shuffling:meiosis}{méiose} à l'autre~; le
nombre de \omterm{def:g10:universal-dna:information}{chromosomes} distincts, donc de gamètes, n'a aucune limite
pratique.

\textbf{18.} $2^{23} \times 2^{23} \approx \num{7e13}$~: sept cents
fois plus de combinaisons que le nombre d'humains ayant jamais vécu, et
cela avant tout \omterm{prop:g12:meiosis-genetic-shuffling:crossingover}{crossing-over}.

\textbf{19.} Les frères et sœurs tirent leurs \omterm{def:g10:universal-dna:gene}{allèles} des mêmes quatre
jeux (deux par parent) et en partagent, en moyenne, la moitié~; des
inconnus tirent dans des jeux différents.

\textbf{20.} 17~\%, la fraction des gamètes de la femelle recombinés
entre $b$ et $vg$, mesure la distance entre les \omterm{def:g10:universal-dna:gene}{gènes}~; le brassage
indépendant des paires et le \omterm{prop:g12:meiosis-genetic-shuffling:crossingover}{crossing-over} à l'intérieur de chaque
paire rendent chaque gamète différent.
\end{solution}
